\vfill\pagebreak
\section{Taking values apart with \texttt{match}}
\label{sec:collections:match}

The \token{match} of Section~\ref{sec:control:match} compared integers with
values and ranges. A pattern can also have the shape of a tuple, of an option or
of a slice. It then takes the value apart, and tests each part against a pattern
of its own.

\subsection{Tuple patterns}

Exercise~1 of Chapter~\ref{chap:control} printed the numbers from 1 to 15,
replacing the multiples of 3 with \textit{Fizz}, those of 5 with \textit{Buzz},
and those of both with \textit{FizzBuzz}. A chain of \token{else if} has to test
the multiples of both first, and computes each remainder twice. A
\token{match} on the pair of remainders lists the four cases as a table.

\lstinputlisting[style=coloredverbatim, caption={FizzBuzz with a tuple pattern, \textit{examples/collections/fizzbuzz.yr}}, label=lst:collections:fizzbuzz]{examples/collections/fizzbuzz.yr}
\vspace{-10pt}%
\begin{lstlisting}[style=bashVerb, escapechar=@+]
1
2
Fizz
4
Buzz
@+\textit{(8 more lines)}@+
14
FizzBuzz
\end{lstlisting}

Line~5 makes a tuple of the two remainders, and matches it. The pattern
\token{(0, 0)} fits a pair whose two fields are 0. In \token{(0, \_)}, the
first field must be 0 and the wildcard accepts any second field, so the arm
catches the multiples of 3 that are not multiples of 5, the pair \token{(0, 0)}
having been caught by the arm above. The arms are tried from top to bottom, as
always, and the order of the cases is the order of the arms.

\subsection{The patterns of the fields}

Each field of a tuple pattern is a pattern of its own, with every form that
Section~\ref{sec:control:match} presented. The next listing uses them all, on a
date written as a pair, the day and then the month.

\lstinputlisting[style=coloredverbatim, caption={The patterns of the fields, \textit{examples/collections/dates.yr}}, label=lst:collections:dates]{examples/collections/dates.yr}
\vspace{-10pt}%
\begin{lstlisting}[style=bashVerb]
(25, 12): Christmas
(1, 1): New Year's Day
(14, 7): a summer day
(3, 11): a day of the first week of a month
(11, 11): a day whose number is that of its month
(20, 11): an ordinary day
\end{lstlisting}

\begin{itemize}
\item \textit{A value.} \token{(25, 12)}, on line~5, fits only the pair whose
  fields are 25 and 12.
\item \textit{Several values.} In \token{(\_, 6 | 7 | 8)}, on line~7, the second
  field must be 6, 7 or 8, and the first can be anything: every day of June,
  July and August.
\item \textit{A range.} In \token{(1 ... 7, \_)}, on line~8, the first field
  must be between 1 and 7.
\item \textit{A name.} \token{(d, m)}, on line~9, fits any pair, and declares
  two variables, \token{d} and \token{m}, that hold its fields. The guard after
  the pattern can then use them, here to compare the day with the month.
\item \textit{The wildcard.} \token{\_}, on line~10, fits any date left.
\end{itemize}

The arms are tried in order. \token{(3, 11)} is in the first week of a month,
and the arm of line~8 catches it before the guard of line~9 is tried.
\token{(1, 1)} fits line~6, and never reaches line~8 either.

\token{main} goes through an array of dates. Each date is a tuple, and
\token{expand date} passes its two fields as the two arguments of
\token{describe} (Section~\ref{sec:collections:expand}).

\subsection{Braces around the arms}

Every arm of \token{describe} is between braces. An arm whose value is not in a
block ends where the value ends, and the next arm starts right after it. When
that next arm starts with a parenthesis, as a tuple pattern does, the compiler
reads the parenthesis as a call of the value before it: without braces, the
first two arms would read as the call \token{"Christmas"(1, 1)}, which makes no
sense. The compiler stops with \errmsg{E2009}{read (, but expected \{}, and
points at the start of the function, far from the fault.

Only the arms followed by a parenthesis need the braces. By convention, though,
a \token{match} has braces on all its arms or on none: when one needs them, all
get them. The next listing places a point of the plane, whose coordinates are
drawn at random, and follows the same rule.

\lstinputlisting[style=coloredverbatim, caption={Placing a point, \textit{examples/collections/point.yr}}, label=lst:collections:point]{examples/collections/point.yr}
\vspace{-10pt}%
\begin{lstlisting}[style=bashVerb, escapechar=@+]
@+\textcolor{teal!80}{alice@dev:\mtilde}@+$ ./point
(-2, 2) is on a diagonal.
@+\textcolor{teal!80}{alice@dev:\mtilde}@+$ ./point
(0, 1) is on the vertical axis.
\end{lstlisting}

The arms of Listing~\ref{lst:collections:fizzbuzz} need no braces: each is a
statement ended by its semicolon, and a semicolon cannot be read as the start of
a call.

\subsection{Option patterns}

\token{Ok(p)} fits an option that holds a value, when that value fits the
pattern \token{p}. With a name in the parentheses, as in
Section~\ref{sec:collections:options}, it fits every value. With another
pattern, it chooses among them.

\lstinputlisting[style=coloredverbatim, caption={Option patterns, \textit{examples/collections/rank.yr}}, label=lst:collections:rank]{examples/collections/rank.yr}
\vspace{-10pt}%
\begin{lstlisting}[style=bashVerb, escapechar=@+]
@+\textcolor{teal!80}{alice@dev:\mtilde}@+$ ./rank
Number: 2
2 is the first prime.
@+\textcolor{teal!80}{alice@dev:\mtilde}@+$ ./rank
Number: 19
19 is the last prime.
@+\textcolor{teal!80}{alice@dev:\mtilde}@+$ ./rank
Number: 7
7 is prime number 4.
\end{lstlisting}

\token{Ok(0)}, on line~14, fits only an option holding 0. On line~15,
\token{Ok(i)} with a guard fits only the last index. Line~16 takes every other
value, and line~17 the empty option.

\subsection{Patterns inside patterns}

Patterns nest: the pattern inside \token{Ok} can be a tuple pattern, whose
fields are patterns in turn. The next listing looks for the best student of a
class, and says whether the grade is perfect.

\lstinputlisting[style=coloredverbatim, caption={Nested patterns, \textit{examples/collections/podium.yr}}, label=lst:collections:podium]{examples/collections/podium.yr}
\vspace{-10pt}%
\begin{lstlisting}[style=bashVerb]
Bruno has a perfect grade.
Emma comes first, with 17.
The class is empty.
\end{lstlisting}

\token{best} returns an option holding a pair, of type
\token{([c8], i32)?}. On line~16, \token{Ok((name, 20))} fits an option that
holds a pair whose grade is 20: the outer parentheses belong to \token{Ok}, the
inner ones to the tuple, and \token{name} holds the first field. Line~17 fits
every other student, and line~18 the empty class.

The same patterns work after \token{if let}. With the \token{point} of
Listing~\ref{lst:collections:point}, \token{if let (0, y) = point \{ ... \}}
runs its block only when the point is on the vertical axis, with \token{y}
holding its second coordinate.

\subsection{Slice patterns}

A slice pattern is written between square brackets, one pattern per element.
It fits a slice that has exactly as many elements, each fitting its pattern. The
last pattern may be followed by \token{...}: it then fits all the remaining
elements, however many, collected into a slice.

\lstinputlisting[style=coloredverbatim, caption={Slice patterns, \textit{examples/collections/sequence.yr}}, label=lst:collections:sequence]{examples/collections/sequence.yr}
\vspace{-10pt}%
\begin{lstlisting}[style=bashVerb]
no value
starts with 0
two values, 4 and 5
1, then [2, 3, 4]
9, then []
\end{lstlisting}

\begin{itemize}
\item \token{[]}, on line~5, fits only the empty slice.
\item \token{[0, \_...]}, on line~6, fits a slice whose first element is 0,
  followed by any number of elements, which the wildcard ignores.
\item \token{[a, b]}, on line~7, fits a slice of exactly two elements, and names
  them.
\item \token{[first, rest...]}, on line~8, fits a slice of at least one
  element. \token{first} holds the first, and \token{rest} the slice of the
  others, which may be empty, as for \token{[9]}.
\end{itemize}

Section~\ref{sec:flow:pattern_matching} lists every pattern of the language.
